
La tabla conserva además las realizaciones cuánticas que su clasificación
abrevia. Los tres leptones cargados y los seis sabores quark pertenecen a la
completación exterior: CAR conduce a Pauli y a Fermi--Dirac, sin atribuir esas
relaciones a un TRIT aislado. Los quarks ocupan treinta y seis celdas al
conservar color y antipartícula antes del cociente. El soporte neutro común
\(\mathcal H_\nu=\bigoplus_{j=1}^{4}\mathcal H_{\nu,G_j}\) tiene dimensiones
\(2+4+6+8=20\). Los símbolos \(G_j\) designan estratos de profundidad del atlas;
\(f_e,f_\mu,f_\tau\) designan la base de interacción. En particular, \(G_4\) es
un estrato del soporte neutro común y no se identifica con un estado estéril o
exótico sin un proyector adicional. PMNS transporta la base de interacción a la
base de autovectores; no crea autovalores, escala de masa ni selector de
profundidad. El fotón conserva ida y retorno dentro de la
carta gauge nula: \(m=0\) no significa evaluar \(\chi_0\) con un exponente muy
negativo. \(W^\pm\) y \(Z^0\) son representaciones bosónicas masivas, con
conjugación de carga para \(W^+\leftrightarrow W^-\). Por último, la
antisección se obtiene mediante inversión y reversión de banda; la masa es par
bajo esa involución y no se duplica como un segundo parámetro independiente.

Todos los sectores conservan el prefijo causal
\[
\text{extensión superviviente}\longrightarrow\text{ruta observable}
\longrightarrow\text{registro cronológico enriquecido}.
\]
Después cambia el lector. En el sector material positivo continúa por
\(\text{firma}\to\text{carácter}\to\text{torres}\to\text{carta dimensional}\).
El sector nulo evita el carácter positivo y fija \(m=0\); los hadrones componen
primero registros cronológicos enriquecidos; los sectores topológicos usan invariantes de
fusión o trenza; y la virtualidad se reconoce por obstrucción terminal. El
dominio de estos mapas es la clasificación completa anterior. Un ensayo
numérico sobre un subconjunto sólo establece el caso finito ensayado y no
reduce ese dominio.

Los censos que acompañan esta tabla pertenecen a cocientes diferentes:
\[
81=25_{\ell^-}+20_\nu+20_d+16_u,
\]
\[
56\ \text{familias con fibras}\quad
41\times1+10\times2+2\times3+1\times4+2\times5=81,
\]
\[
13\ \text{multisecciones},\qquad324=81\times4\ \text{genealogías}.
\]
Asimismo, las \(192\)-A clases de representación asociadas a
\(\operatorname{diag}(2,6,2,2,4)\) y el índice reticular \(192\)-B de
\(\operatorname{SNF}(2,2,2,2,12)\), con \(52\) firmas crudas, no son el
El cardinal del catálogo expositivo histórico no cuenta partículas individuales ni define la taxonomía material.
no partículas individuales.
